\section*{Chapitre \ref{ch:b2:orbital-forcing} --- Forçage orbital et âges glaciaires}

\begin{solution}{exo:b2:orbital-forcing:1}
(a) $1.0167/0.9833 = 1.0340$~: l'aphélie est \qty{3.4}{\percent} plus loin que le périhélie.
(b) $a(1 + e) - a(1 - e) = 2ea = 0.0334\,a$, soit \qty{3.34}{\percent} de $a$, proche des 3.4 pour cent
généralement cités. (c) Le rayonnement solaire décroît comme le carré inverse de la distance : $1.0340^2 = 1.069$, soit
environ \qty{7}{\percent} de rayonnement en plus au périhélie.
\end{solution}

\begin{solution}{exo:b2:orbital-forcing:2}
(a) Aujourd'hui $0.017236 \times \sin\ang{101.37} = 0.017236 \times 0.9804 = 0.01690$~; il y a \num{10000} ans
$0.019971 \times \sin\ang{293.86} = 0.019971 \times (-0.9145) = -0.01826$~; il y a \num{6000} ans
$0.019249 \times \sin\ang{359.99} \approx 0$, valeurs des tables. (b) La longitude du Soleil au
périhélie est $\varpi + \ang{180}$~: \ang{281.37} aujourd'hui, \ang{473.86} $-$ \ang{360} $=$ \ang{113.86}
il y a \num{10000} ans, \ang{179.99} il y a \num{6000} ans. (c) Il y a \num{10000} ans, le périhélie
tombait \ang{24} après le solstice de juin, le plus proche des trois de celui-ci, pendant l'été boréal~;
il y a \num{6000} ans, il tombait à l'équinoxe de septembre, où la \omterm{def:b2:orbital-forcing:elements}{précession climatique} est nulle.
\end{solution}

\begin{solution}{exo:b2:orbital-forcing:3}
(a) Les fréquences s'ajoutent : $1/25.77 + 1/112 = 0.03880 + 0.00893 = \qty{0.04773}{ka^{-1}}$, soit une période
de $1/0.04773 = 21.0$ milliers d'années. (b) Le résultat tombe dans l'intervalle des périodes principales,
de 19.0 à 23.7 milliers d'années. Un seul nombre ne peut pas les reproduire : la rotation de l'orbite
est elle-même la somme de plusieurs mouvements de fréquences différentes, dont chacun se combine avec la
précession de l'axe pour donner une période de la \omterm{def:b2:orbital-forcing:elements}{précession climatique}~; la moyenne est la valeur de
\num{23000} ans généralement citée.
\end{solution}

\begin{solution}{exo:b2:orbital-forcing:4}
$\sqrt{1 - 0.0167^2} = 0.99986$. Pour \ang{22.1}~: $1360 \times 0.37622/(\pi \times 0.99986) =
\qty{162.9}{W.m^{-2}}$~; pour \ang{24.5}~: $1360 \times 0.41469/(\pi \times 0.99986) =
\qty{179.5}{W.m^{-2}}$. Le rapport est $\sin\ang{24.5}/\sin\ang{22.1} = 1.102$~: la moyenne annuelle au
pôle croît de \qty{10.2}{\percent}, soit environ \qty{16.7}{W.m^{-2}}, entre les deux extrêmes.
\end{solution}

\begin{solution}{exo:b2:orbital-forcing:5}
(a) Au solstice de juin $(a/r)^2 = \bigl((1 - e\sin\varpi)/(1 - e^2)\bigr)^2$. Aujourd'hui $e\sin\varpi
= 0.016898$ et $(0.983102/0.999703)^2 = 0.9671$, contre $1 - 0.033796 = 0.9662$~;
il y a \num{116000} ans $e\sin\varpi = 0.043939$ et $(0.956061/0.998065)^2 = 0.9176$, contre $1 -
0.087878 = 0.9121$. (b) Au solstice de décembre $(a/r)^2 = \bigl((1 + e\sin\varpi)/(1 -
e^2)\bigr)^2$~: aujourd'hui $1.0347$ contre $1.0338$, il y a \num{116000} ans $1.0940$ contre $1.0879$.
(c) Les termes négligés sont d'ordre $e^2$~: $(1 - x)^2/(1 - e^2)^2 \approx 1 - 2x + x^2 + 2e^2$
avec $x = e\sin\varpi$, ici environ $3e^2$. Ils valent $3 \times 0.0172^2 = 0.0009$ aujourd'hui et $3
\times 0.0440^2 = 0.0058$ il y a \num{116000} ans, proches des écarts trouvés, $0.0009$ et
$0.0055$~: l'approximation est six fois moins bonne lorsque l'excentricité est deux fois et demie
plus grande.
\end{solution}

\begin{solution}{exo:b2:orbital-forcing:6}
Au solstice de juin $\delta = \varepsilon = \ang{23.446}$. $\cos h_0 = -\tan\ang{65}\tan\ang{23.446}
= -2.14451 \times 0.43369 = -0.93006$, donc $h_0 = \ang{158.44} = \qty{2.7654}{rad}$. Le jour dure
$2 \times 158.44/15 = 21.13$~heures (la Terre tourne de \ang{15} par heure). Avec $\sin\ang{65} = 0.90631$,
$\cos\ang{65} = 0.42262$, $\sin\delta = 0.39788$, $\cos\delta = 0.91744$ et $\sin h_0 = 0.36742$,
le facteur entre parenthèses vaut $2.7654 \times 0.90631 \times 0.39788 + 0.42262 \times 0.91744 \times 0.36742 = 0.99720
+ 0.14246 = 1.13966$. À la distance moyenne $\bar Q = (1360/\pi) \times 1.13966 = 432.90 \times 1.13966
= \qty{493.4}{W.m^{-2}}$~; avec le facteur de distance de l'\cref{exo:b2:orbital-forcing:5}, $493.4 \times
0.9671 = \qty{477.1}{W.m^{-2}}$. Pour vérification, les mêmes étapes à \ang{60}~N donnent \qty{475.96}{W.m^{-2}}, là où les tables donnent \num{475.95}.
\end{solution}

\begin{solution}{exo:b2:orbital-forcing:7}
(a) L'équinoxe de mars est $\lambda = 0$, celui de septembre $\lambda = \pi$. Avec $\sqrt{1 - e^2}
\approx 1$, la durée vaut $\frac{T}{2\pi}\int_0^{\pi}\bigl(1 + 2e\cos(\lambda -
\varpi)\bigr)\,\mathrm d\lambda = \frac{T}{2\pi}\bigl(\pi + 2e[\sin(\lambda - \varpi)]_0^{\pi}\bigr)
= \frac{T}{2\pi}(\pi + 4e\sin\varpi) = \frac T2 + \frac{2T}{\pi}e\sin\varpi$. (b) L'autre demi-année
est plus courte de la même quantité, donc la différence vaut $(4T/\pi)\,e\sin\varpi = (4 \times
365.2422/\pi) \times 0.016898 = 465.04 \times 0.016898 = 7.86$~jours. (c) Le printemps et l'été ensemble
dépassent l'automne et l'hiver de $4.5 + 3 = 7.5$~jours~; les deux écarts cités sont arrondis, et l'accord
se situe dans cet arrondi.
\end{solution}

\begin{solution}{exo:b2:orbital-forcing:8}
(a) Pour $e = 0.0046$~: $\sqrt{1 - e^2} = 0.999989$ et $340/0.999989 = \qty{340.004}{W.m^{-2}}$. Pour
$e = 0.0565$~: $\sqrt{1 - 0.003192} = 0.998403$ et $340/0.998403 = \qty{340.544}{W.m^{-2}}$. (b) La
différence vaut \qty{0.54}{W.m^{-2}}, soit \qty{0.16}{\percent}. (c) L'insolation de mi-juillet à \ang{65}~N
a varié de $497 - 388 = \qty{109}{W.m^{-2}}$, soit environ \qty{25}{\percent} de sa moyenne et deux cents
fois la variation de la moyenne annuelle globale : l'orbite agit sur le climat en déplaçant le rayonnement entre
les saisons et les latitudes, non en modifiant le total.
\end{solution}

\begin{solution}{exo:b2:orbital-forcing:9}
Avec $\omega = 2\pi/P$ et $\tau = \qty{10}{ka}$~: pour $P = 23$, $\omega\tau = 2.732$, une amplitude
relative $1/\sqrt{1 + 7.464} = 0.344$ et un retard $\arctan(2.732)/\omega = 1.2199/0.2732 = 4.47$~ka~;
pour $P = 41$, $\omega\tau = 1.532$, amplitude $0.546$, retard $0.9926/0.1532 = 6.48$~ka~; pour $P =
100$, $\omega\tau = 0.628$, amplitude $0.847$, retard $0.5610/0.0628 = 8.93$~ka. La calotte glaciaire
favorise le forçage lent : elle suit \qty{85}{\percent} d'un cycle de \num{100000} ans mais seulement un
tiers d'un cycle de précession, car elle n'a pas le temps de croître ou de décroître avant qu'un forçage rapide
ne s'inverse. Elle agit comme un filtre passe-bas.
\end{solution}

\begin{solution}{exo:b2:orbital-forcing:10}
(a) L'équation est linéaire : si $V_1$ répond à $A_1\sin\omega_1t$ et $V_2$ à
$A_2\sin\omega_2t$, $V_1 + V_2$ répond à leur somme. Par la \cref{prop:b2:orbital-forcing:lag}, $V_1$ et
$V_2$ sont des sinusoïdes de fréquences $\omega_1$ et $\omega_2$ (une fois le transitoire éteint), si bien que
la réponse ne contient aucune autre fréquence. (b) Avec $V \approx B_1\sin(\omega_1t - \phi_1) +
B_2\sin(\omega_2t - \phi_2)$,
\[
\begin{aligned}
  V^2 &= \tfrac12(B_1^2 + B_2^2) - \tfrac12B_1^2\cos(2\omega_1t - 2\phi_1) -
  \tfrac12B_2^2\cos(2\omega_2t - 2\phi_2)\\
  &\quad + B_1B_2\bigl[\cos\bigl((\omega_1 - \omega_2)t - \phi_1 + \phi_2\bigr) -
  \cos\bigl((\omega_1 + \omega_2)t - \phi_1 - \phi_2\bigr)\bigr],
\end{aligned}
\]
donc le terme $kV^2$ force l'équation aux fréquences de différence et de somme, et aux fréquences
doubles. (c) $1/(1/19.0 - 1/23.7) = 1/(0.05263 - 0.04219) = 95.8$ milliers d'années~; la paire 19.0
et 22.4 donne 125 milliers d'années. Les fréquences de différence des termes de précession tombent dans
la bande de l'excentricité, comme il se doit, puisque la \omterm{def:b2:orbital-forcing:elements}{précession climatique} est $e\sin\varpi$. Une calotte glaciaire
non linéaire pourrait donc tirer un rythme de \num{100000} ans de la précession~; une calotte linéaire ne le peut pas.
Les cycles de \num{100000} ans sont la preuve d'un comportement non linéaire, bien que ce calcul seul
ne dise pas lequel.
\end{solution}

\begin{solution}{exo:b2:orbital-forcing:11}
(a) Comme dans la démonstration de la \cref{prop:b2:orbital-forcing:annual}, $(a/r)^2\,\mathrm dt =
T\,\mathrm d\lambda/(2\pi\sqrt{1 - e^2})$, et $\sin\delta = \sin\varepsilon\sin\lambda$ est positif
pour $0 < \lambda < \pi$. Donc $E = \frac{S_0T\sin\varepsilon}{2\pi\sqrt{1 - e^2}}\int_0^{\pi}\sin\lambda
\,\mathrm d\lambda = S_0T\sin\varepsilon/(\pi\sqrt{1 - e^2})$. (b) $T = 365.25 \times 86400 =
\qty{3.1558e7}{s}$, $S_0T/\pi = \qty{1.3661e10}{J.m^{-2}}$~; pour \ang{22.1}, $E = 1.3661 \times 10^{10}
\times 0.37622/0.99986 = \qty{5.14}{GJ.m^{-2}}$~; pour \ang{24.5}, $\qty{5.67}{GJ.m^{-2}}$,
soit \qty{10}{\percent} de plus. (c) La précession fixe la façon dont la même énergie se répartit dans le temps : un été
au périhélie est plus court et plus intense, un été à l'aphélie plus long et plus faible. La fonte ne commence
qu'au-dessus de \qty{0}{\celsius}, si bien qu'un été court et intense peut faire fondre davantage qu'un été long et frais
recevant le même total.
\end{solution}

\begin{solution}{exo:b2:orbital-forcing:12}
(a) $3.93\log_2(279.4/189.7) = 3.93\log_2 1.4729 = 3.93 \times 0.5586 = \qty{2.20}{W.m^{-2}}$. (b) La
plus petite hausse va de 194 à \qty{266}{ppm}~: $3.93\log_2 1.3711 = 3.93 \times 0.4554 =
\qty{1.79}{W.m^{-2}}$~; la plus grande, de 188 à \qty{282}{ppm}~: $3.93\log_2 1.5 = 3.93 \times 0.5850 =
\qty{2.30}{W.m^{-2}}$. (c) Le forçage du dioxyde de carbone vaut environ quatre fois la plus grande variation de
l'insolation annuelle globale, \qty{0.54}{W.m^{-2}}, et il agit partout et toute l'année~; l'insolation estivale à \ang{65}~N
varie d'environ \qty{109}{W.m^{-2}}, mais en une seule latitude et une seule saison. L'orbite est le déclencheur,
par les étés des calottes glaciaires nord~; le dioxyde de carbone, avec la \omterm{def:b2:precambrian-earth:ice-albedo}{rétroaction glace--albédo}, est un amplificateur
qui étend le changement à tout le globe.
\end{solution}

\begin{solution}{pb:b2:orbital-forcing:1}
\textbf{1.} Alors $0.0440 \times \sin\ang{92.7} = 0.0440 \times 0.99889 = 0.04395$~; aujourd'hui $0.0172 \times
\sin\ang{101.4} = 0.0172 \times 0.98027 = 0.01686$, environ 2.6 fois moins.
\textbf{2.} Alors $1 - 0.0440 = 0.956$ et $1 + 0.0440 = 1.044$~; aujourd'hui $0.9828$ et $1.0172$.
\textbf{3.} En $\lambda = \varpi + \ang{180}$~: alors \ang{272.7}, aujourd'hui \ang{281.4}, toutes deux proches du
solstice de décembre ($\lambda = \ang{270}$)~; alors, le périhélie tombait moins de trois jours après ce solstice,
et l'aphélie au solstice de juin.
\textbf{4.} Alors $(1.044/0.956)^2 = 1.0921^2 = 1.193$, aujourd'hui $(1.0172/0.9828)^2 = 1.0350^2 = 1.071$~:
\qty{19}{\percent} de rayonnement en plus au périhélie alors, \qty{7}{\percent} aujourd'hui.
\textbf{5.} Plus petite, \ang{22.52} contre \ang{23.45}~: le Soleil d'été était plus bas aux hautes latitudes,
et les hautes latitudes recevaient moins de rayonnement en été et sur l'année.
\textbf{6.} Alors $340/\sqrt{1 - 0.001936} = 340/0.999032 = \qty{340.33}{W.m^{-2}}$~; aujourd'hui $340/0.999852
= \qty{340.05}{W.m^{-2}}$. L'écart, \qty{0.28}{W.m^{-2}} ou moins d'un millième,
est négligeable.
\textbf{7.} En $\lambda = \ang{90}$, $\cos(\ang{90} - \varpi) = \sin\varpi$ dans
\cref{prop:b2:orbital-forcing:distance}.
\textbf{8.} Alors $(1 - 0.001936)/(1 - 0.04395) = 0.998064/0.956049 = 1.04395$~; aujourd'hui $(1 -
0.000296)/(1 - 0.01686) = 0.999704/0.983139 = 1.01685$.
\textbf{9.} Alors $1/1.04395^2 = 0.91758$~; aujourd'hui $1/1.01685^2 = 0.96714$.
\textbf{10.} Alors $1 - 2 \times 0.04395 = 0.91210$~; aujourd'hui $1 - 2 \times 0.01686 = 0.96628$~: l'approximation
se trompe de $0.0055$ alors, de $0.0009$ aujourd'hui.
\textbf{11.} Alors $1 - 0.91758 = \qty{8.2}{\percent}$ plus faible, aujourd'hui \qty{3.3}{\percent}.
\textbf{12.} Au solstice $\delta = \varepsilon$~: \ang{22.52} alors, \ang{23.45} aujourd'hui.
\textbf{13.} Alors $\cos h_0 = -2.14451 \times 0.41462 = -0.88916$, $h_0 = \ang{152.77} =
\qty{2.6663}{rad}$~; aujourd'hui $\cos h_0 = -2.14451 \times 0.43378 = -0.93023$, $h_0 = \ang{158.47} =
\qty{2.7659}{rad}$.
\textbf{14.} $2h_0/\ang{15}$ par heure : alors $20.37$~heures, aujourd'hui $21.13$~heures.
\textbf{15.} Alors $2.6663 \times 0.90631 \times 0.38301 + 0.42262 \times 0.92375 \times 0.45759 =
0.92554 + 0.17864 = 1.10418$ (\num{1.10417} sans arrondi), et $432.90 \times 1.10417 =
\qty{478.0}{W.m^{-2}}$~; aujourd'hui $2.7659 \times 0.90631 \times 0.39795 + 0.42262 \times 0.91741 \times
0.36697 = 0.99757 + 0.14228 = 1.13985$ (\num{1.13982} sans arrondi), et $432.90 \times 1.13982 = \qty{493.4}{W.m^{-2}}$.
\textbf{16.} Alors $478.0 \times 0.91758 = \qty{438.6}{W.m^{-2}}$~; aujourd'hui $493.4 \times 0.96714 =
\qty{477.2}{W.m^{-2}}$, soit \qty{38.6}{W.m^{-2}} de plus qu'alors.
\textbf{17.} À mi-juillet, la longitude du Soleil vaut \ang{120}~: $\sin\delta = \sin\varepsilon\sin\ang{120}$,
une déclinaison d'environ \ang{20} au lieu de \ang{23}, donc un jour plus court et un Soleil plus bas~: à
\ang{65}~N, l'insolation journalière moyenne est maximale près du solstice, lorsque la déclinaison est la plus grande.
\textbf{18.} Avec la distance d'aujourd'hui, le seul changement d'obliquité coûte $(493.4 - 478.0) \times
0.96714 = \qty{14.9}{W.m^{-2}}$~; avec l'ancienne obliquité, la seule distance coûte $478.0 \times
(0.96714 - 0.91758) = \qty{23.7}{W.m^{-2}}$. Les deux s'additionnent à \qty{38.6}{W.m^{-2}}~; la distance
compte pour environ trois cinquièmes.
\textbf{19.} Après un changement brusque, la calotte glaciaire a atteint $1 - \mathrm e^{-t/\tau}$ de celui-ci~:
$1 - \mathrm e^{-0.5} = 0.39$ après \qty{5}{ka}, $1 - \mathrm e^{-1} = 0.63$ après \qty{10}{ka},
$1 - \mathrm e^{-2} = 0.86$ après \qty{20}{ka}.
\textbf{20.} $\omega\tau = 2\pi \times 10/23 = 2.732$~: la calotte glaciaire suit $1/\sqrt{1 + 2.732^2} =
0.34$ du forçage, avec un retard de $\arctan(2.732)/\omega = 4.5$~ka.
\textbf{21.} Une demi-période de précession, soit environ 11 à \qty{12}{ka}, est à peu près égale à un temps de réponse~: avant
que la calotte glaciaire ait couvert les deux tiers du chemin, l'insolation s'inverse et la ramène en arrière.
\textbf{22.} Dans le modèle linéaire, la glace croît et décroît au même rythme pour des forçages de même
taille~: un forçage croissant et un forçage décroissant donnent des réponses symétriques. Une accumulation lente suivie
d'un effondrement rapide exige une réponse qui dépend de l'état ou du sens du changement, c'est-à-dire
une réponse non linéaire.
\textbf{23.} Le dioxyde de carbone a augmenté de $3.93\log_2(266/194) = 1.8$ à $3.93\log_2(282/188) =
\qty{2.3}{W.m^{-2}}$ du glaciaire à l'\omterm{def:b2:orbital-forcing:glacials}{interglaciaire} : un forçage sur tout le globe et pendant toute l'année.
Libéré lorsque la circulation atlantique s'affaiblit, il porte le changement des étés boréaux à toute la planète et,
avec l'albédo de la glace qui se réduit, l'amplifie.
\textbf{24.} Lorsque le dernier \omterm{def:b2:orbital-forcing:glacials}{interglaciaire} s'est achevé, l'insolation journalière moyenne du solstice de juin à
\ang{65}~N était d'environ \textbf{\qty{439}{W.m^{-2}}}, \qty{39}{W.m^{-2}} de moins qu'aujourd'hui et sa plus faible
valeur du cycle.
\end{solution}
